\chapter{Pipeline Implementation}
\label{cha:implementation}

This chapter presents the software implementation of the signal denoising~(ATP-1)
and data compression~(ATP-3) stages of the pre-processing pipeline. The
implementation is structured as a set of independent Python modules, each
encapsulating one pipeline stage so that any stage can be replaced or extended
without modifying the others. All code is written in Python~3.10 and is publicly
available in the project repository at
\url{https://github.com/Sravya0304/pyrometer-preprocessing}.

\section{Software Architecture}
\label{sec:architecture}

The pipeline is organised into four layers that communicate through NumPy float32
arrays passed directly in memory. No intermediate files are written between
stages during a full evaluation run.

The \textbf{data loading layer} is implemented in \texttt{layer01\_pipeline.py}.
It reads the NIST \texttt{Layer01.mat} file using \texttt{scipy.io.loadmat},
extracts the per-frame maximum pixel value from the three-dimensional radiance
array of shape $126 \times 360 \times N_\text{frames}$, and converts it to
degrees Celsius using the Sakuma--Hattori calibration equation:
\begin{equation}
T\,(\text{\textdegree C}) = A \cdot x_\text{raw} + B - 273.15
\label{eq:sh-impl}
\end{equation}
where $A = 2.655$ and $B = -800.7$ for Layer~01. All frames with
$T \leq 10$\,\textdegree C are discarded as inactive, leaving~2,065 non-zero
frames.

The \textbf{denoising layer} contains three scripts. \texttt{denoise.py} exports
the primary \texttt{denoise\_signal()} function used by the main pipeline.
\texttt{classical\_methods.py} implements all five classical filters with
consistent function signatures. \texttt{ml\_denoise.py} implements the four ML
denoisers, each exposing a \texttt{train()} and \texttt{predict()} interface.

The \textbf{compression layer} mirrors the denoising layer. \texttt{compress.py}
exports wavelet and SVD compression and reconstruction functions.
\texttt{ml\_compress.py} implements PCA and Autoencoder compression.

The \textbf{comparison and visualisation layer} consists of
\texttt{d3\_comparison.py}, which runs all methods under identical conditions and
records metrics to CSV files; \texttt{visualise.py}, which produces the D4
dashboard; and \texttt{d5\_analysis.py}, which evaluates all methods across all
ten NIST layers.

\section{Data Pre-processing}
\label{sec:preprocessing}

Before any denoising or compression is applied, the raw signal is passed through
three consistent pre-processing steps.

\textbf{Outlier clipping.} Temperature values below~$0$\,\textdegree C or above
$3000$\,\textdegree C are clipped to the boundary. These values are physically
impossible for an in-process pyrometer measurement.

\textbf{Zero-frame removal.} Frames at or below~$10$\,\textdegree C are removed.
These correspond to laser-off periods where the melt pool has solidified.
Including them would bias noise statistics and confuse ML training.

\textbf{Simulated thermocouple reference.} Because the NIST dataset contains no
physical thermocouple, a reference signal is generated by applying a first-order
low-pass filter with time constant~$\alpha = 0.08$ to the calibrated signal:
\begin{equation}
T_\text{TC}[i] = T_\text{TC}[i-1] +
  \alpha\bigl(T[i] - T_\text{TC}[i-1]\bigr)
\label{eq:lpf-impl}
\end{equation}
This reference simulates thermocouple thermal lag and serves as the clean
training target for all four ML denoisers.

\section{Denoising Implementation (ATP-1)}
\label{sec:impl-denoise}

\subsection{Classical Filters}

The \textbf{moving average} is implemented with \texttt{numpy.convolve} and a
uniform kernel of width~$k = 7$. Edge effects are handled by the \texttt{same}
padding mode. The kernel width was chosen after testing $k \in \{3, 5, 7, 11\}$;
smaller values left excess residual noise while larger values distorted the peak
temperature frame.

The \textbf{median filter} uses \texttt{scipy.signal.medfilt} with kernel
size~7. The median is insensitive to extreme values, so a single spike is replaced
by the middle value of its neighbours without affecting surrounding samples. The
median filter is used as the primary denoiser in \texttt{denoise.py} and also as
the clean training target for all four ML denoisers.

The \textbf{Savitzky--Golay filter} uses \texttt{scipy.signal.savgol\_filter}
with window~$N=11$ and polynomial degree~$p=3$~\cite{savitzky1964smoothing}. The
cubic polynomial can follow the curvature of a genuine temperature peak,
preserving peak heights and widths better than the moving average while still
suppressing high-frequency noise.

The \textbf{Gaussian filter} uses \texttt{scipy.ndimage.gaussian\_filter1d} with
$\sigma = 3.0$. The Gaussian kernel weights central samples more heavily than
edge samples, producing less spectral distortion than the moving average at
comparable noise suppression.

The \textbf{Kalman filter} is a recursive loop with $Q = 0.001$ and $R = 10.0$.
At each time step the gain and state estimate are updated:
\begin{align}
P_i^-     &= P_{i-1} + Q \\
K_i       &= \frac{P_i^-}{P_i^- + R} \\
\hat{x}_i &= \hat{x}_{i-1} + K_i\,(x_i - \hat{x}_{i-1}) \\
P_i       &= (1 - K_i)\,P_i^-
\end{align}
The initial state is set to the first sample and $P_0 = 1.0$. The ratio
$Q/R = 10^{-4}$ means the filter trusts its model prediction strongly, producing
strong Gaussian noise suppression. It is the only denoising method capable of
real-time processing because it operates one sample at a time without storing
signal history.

\subsection{ML Denoisers}

All four ML denoisers are implemented in \texttt{ml\_denoise.py} using
PyTorch~\cite{pytorch}. The signal is normalised to~$[0,1]$ using
\texttt{MinMaxScaler} and segmented into overlapping windows of $W=32$ samples
with~50\,\% overlap. The median-filtered signal serves as the clean training
target for all models. Training settings for all four: batch size~32, Adam
optimiser~\cite{kingma2014adam} at learning rate~$10^{-3}$, MSE loss.

The \textbf{CNN denoiser} uses four Conv1d layers:
\begin{itemize}[nosep, leftmargin=2em]
\item Conv1d(1$\to$16, kernel=7, padding=3) + ReLU
\item Conv1d(16$\to$32, kernel=5, padding=2) + ReLU
\item Conv1d(32$\to$16, kernel=5, padding=2) + ReLU
\item Conv1d(16$\to$1, kernel=7, padding=3)
\end{itemize}
Symmetric zero-padding preserves the input length at every layer. Trained for
50 epochs.

The \textbf{LSTM denoiser} uses two stacked layers with hidden size~64 followed
by a Linear(64$\to$1) projection. The gating mechanism retains the long-term
temperature trend while suppressing short-duration noise transients. Trained for
50 epochs.

The \textbf{Autoencoder denoiser} uses a fully connected encoder--decoder with
bottleneck size~8: Linear(32$\to$16)$\to$ReLU$\to$Linear(16$\to$8) for the
encoder and the mirror for the decoder. The narrow bottleneck forces noise
removal because noise has high entropy and cannot be efficiently encoded in
8 values. Trained for 50 epochs.

The \textbf{BiLSTM denoiser} extends the LSTM with \texttt{bidirectional=True}.
Forward and backward hidden states are concatenated before the
Linear(128$\to$1) projection, giving each predicted sample access to both past
and future context within the window. Trained for 50 epochs.

Figure~\ref{fig:pipeline} shows the full pipeline output on NIST Layer~01,
including the denoising stage output in the top panel.

\begin{figure}[H]
\centering
\includegraphics[width=\textwidth]{figures/pipeline_result}
\caption{Three-stage pipeline output on NIST Layer~01. Top panel: raw signal
         versus Kalman-denoised output. Middle panel: calibrated signal
         versus the simulated thermocouple reference. Bottom panel: original
         calibrated signal versus wavelet-compressed reconstruction.}
\label{fig:pipeline}
\end{figure}

Figure~\ref{fig:classical} shows all five classical denoising method outputs
on Layer~01. The Kalman filter produces the smoothest output, consistent with
its strong performance on Gaussian noise.

\begin{figure}[H]
\centering
\includegraphics[width=\textwidth]{figures/classical_results}
\caption{Classical denoising method outputs on NIST Layer~01. Top panels show
         the full time series for each method overlaid on the raw signal.
         Bottom panels zoom in around the peak temperature frame to compare
         noise suppression and peak preservation across methods.}
\label{fig:classical}
\end{figure}

Figure~\ref{fig:ml-denoise} shows the four ML denoiser outputs on the test set.

\begin{figure}[H]
\centering
\includegraphics[width=\textwidth]{figures/ml_denoise_result}
\caption{Machine learning denoiser outputs on the NIST Layer~01 test set.
         Each panel overlays the raw signal~(grey), the median-filtered training
         target~(blue), and the ML denoiser prediction~(orange). The CNN and
         LSTM follow the training target closely.}
\label{fig:ml-denoise}
\end{figure}

\section{Compression Implementation (ATP-3)}
\label{sec:impl-compress}

All compression methods receive the calibrated temperature signal from the
calibration stage~\cite{ajay2026} as input.

\textbf{Downsampling} retains every other sample (\texttt{signal[::2]}) and
reconstructs via \texttt{numpy.interp} linear interpolation. Fixed
ratio of~2$\times$.

\textbf{Truncated SVD} reshapes the signal into non-overlapping windows of~32
samples to form a matrix~$\mathbf{X} \in \mathbb{R}^{m \times 32}$, applies
\texttt{numpy.linalg.svd}, and retains the~$k=10$ largest singular
triplets~\cite{golub2013matrix}:
\begin{equation}
\mathbf{X} \approx \mathbf{U}_k\,\boldsymbol{\Sigma}_k\,\mathbf{V}_k^\top
\end{equation}
Achieved ratio is~2.1$\times$ due to float64 storage overhead.

\textbf{Wavelet compression} uses PyWavelets~\cite{pywavelets} with the Haar
transform. The bottom~90\,\% of coefficients by absolute magnitude are zeroed;
only non-zero coefficients and their int32 indices are stored. Achieved ratio:
5.0$\times$ with RMSE~36.55\,\textdegree C --- the only method satisfying
the ATP-3 target.

\textbf{PCA compression} applies \texttt{sklearn.decomposition.PCA(n\_components=3)}
to normalised overlapping windows of $W=32$~\cite{sklearn,jolliffe2002pca}.
Ratio: $32/3 \approx 10.7\times$.

\textbf{Autoencoder compression} uses a three-layer encoder~(32$\to$16$\to$8$\to$3)
and mirror decoder, trained for~100 epochs with Adam and MSE loss. Same
ratio as PCA: 10.7$\times$.

Figure~\ref{fig:ml-compress} shows the reconstruction quality for both ML
compression methods.

\begin{figure}[H]
\centering
\includegraphics[width=\textwidth]{figures/ml_compress_result}
\caption{ML compression reconstruction on NIST Layer~01. Top: PCA~($n=3$,
         ratio~10.7$\times$). Bottom: Autoencoder~(bottleneck=3,
         ratio~10.7$\times$). Original signal in blue; reconstruction in orange.}
\label{fig:ml-compress}
\end{figure}

\section{Comparison Script (D3)}
\label{sec:d3}

\texttt{d3\_comparison.py} runs all ten denoising methods and all six compression
methods on Layer~01 under a single 80/20 temporal split. Results are saved to
three CSV summary files and a six-panel comparison figure. Figure~\ref{fig:d3}
shows this full D3 output.

\begin{figure}[H]
\centering
\includegraphics[width=\textwidth]{figures/d3_comparison}
\caption{D3 classical versus ML comparison on NIST Layer~01 IN625 data. Top row:
         denoising signal overlay and noise standard deviation bar chart.
         Middle row: calibration results from the parallel thesis~\cite{ajay2026}.
         Bottom row: compression overlay and ratio--RMSE scatter plot.}
\label{fig:d3}
\end{figure}

\section{Visualisation Dashboard (D4)}
\label{sec:d4}

\texttt{visualise.py} produces the D4 multi-panel dashboard shown in
Figure~\ref{fig:d4}. The dashboard gives a process engineer a complete overview
of the signal at every pipeline stage, with event markers at the hotspot frame
and laser transition point.

\begin{figure}[H]
\centering
\includegraphics[width=\textwidth]{figures/d4_dashboard}
\caption{D4 visualisation dashboard for NIST Layer~01. Panels show the raw
         signal, denoised output, calibrated output, and compressed
         reconstruction with event markers. Quantitative metrics accompany
         each signal plot to support method selection.}
\label{fig:d4}
\end{figure}

\section{Trade-off Analysis (D5)}
\label{sec:d5}

\texttt{d5\_analysis.py} evaluates all methods across all ten NIST layers and
records per-layer metrics to \texttt{d5\_summary.csv}. The cross-layer evaluation
confirms whether Layer~01 results are representative of the full NIST dataset.
Figure~\ref{fig:d5} shows the D5 analysis output.

\begin{figure}[H]
\centering
\includegraphics[width=\textwidth]{figures/d5_analysis}
\caption{D5 cross-layer trade-off analysis across all ten NIST layers. Bar
         charts show noise reduction and compression RMSE per method. The
         scatter plot shows the compression ratio versus RMSE frontier, with
         the vertical dashed line marking the~4$\times$ ATP-3 target.}
\label{fig:d5}
\end{figure}

\section{Code Organisation}
\label{sec:code-org}

Table~\ref{tab:scripts} lists all fifteen pipeline scripts with their
deliverable mapping and role. All scripts follow consistent function interfaces
and have module-level docstrings. Random seeds are fixed at~42
(\texttt{numpy.random.seed(42)}, \texttt{torch.manual\_seed(42)}) to ensure
reproducibility across ML training runs. The complete codebase is available at
\url{https://github.com/Sravya0304/pyrometer-preprocessing}.

\begin{table}[H]
\centering
\caption{Complete list of pipeline scripts, deliverables, and roles.}
\label{tab:scripts}
\begin{tabular}{lll}
\toprule
\textbf{Script} & \textbf{Deliverable} & \textbf{Role} \\
\midrule
\texttt{layer01\_pipeline.py}        & D1 & NIST data loading and pipeline runner \\
\texttt{two\_pyrometer\_pipeline.py} & D1 & Two-pyrometer simulation \\
\texttt{denoise.py}                  & D2 & Primary denoising module \\
\texttt{calibrate.py}                & D2 & Calibration module \\
\texttt{compress.py}                 & D2 & Compression module \\
\texttt{pipeline\_stages.py}         & D2 & Stage-by-stage pipeline runner \\
\texttt{classical\_methods.py}       & D3 & All classical methods \\
\texttt{ml\_denoise.py}              & D3 & CNN, LSTM, AE, BiLSTM denoisers \\
\texttt{ml\_calibrate.py}            & D3 & RF, MLP, GB, SVR calibrators \\
\texttt{ml\_compress.py}             & D3 & PCA, AE, VAE, Deep AE compressors \\
\texttt{d3\_comparison.py}           & D3 & Classical vs ML full comparison \\
\texttt{visualise.py}                & D4 & Multi-panel dashboard \\
\texttt{d5\_analysis.py}             & D5 & Cross-layer trade-off analysis \\
\texttt{load\_podfam.py}             & --- & Single PODFAM file pipeline \\
\texttt{podfam\_all\_files.py}       & --- & All PODFAM files evaluation \\
\bottomrule
\end{tabular}
\end{table}

\let\cleardoublepage\clearpage
